\section*{अध्याय \ref{ch:g4:plants-without-seeds} --- बिना बीज वाले पौधे}

\begin{solution}{exo:g4:plants-without-seeds:1}
\omterm{prop:g4:plants-without-seeds:damp}{मॉस} और \omterm{prop:g4:plants-without-seeds:damp}{फ़र्न}। वे \omterm{def:g4:plants-without-seeds:spore}{बीजाणुओं} से जनन करते हैं --- यानी बहुत बड़ी संख्या में
छोड़े गए सूक्ष्म कणों से।
\end{solution}

\begin{solution}{exo:g4:plants-without-seeds:2}
\omterm{def:g4:plants-without-seeds:spore}{बीजाणु} सूक्ष्म और नंगा होता है --- न \omterm{def:g3:animals-through-seasons:active}{भोजन-भंडार}, न पहले से बना नन्हा
\omterm{def:g1:plants-around-us:plant}{पौधा}; \omterm{def:g2:from-seed-to-plant:seed}{बीज} दोनों लिए रहता है, और वह भी एक रक्षक छिलके के भीतर। और
\omterm{def:g4:plants-without-seeds:spore}{बीजाणु} लाखों में बनते हैं, \omterm{def:g2:from-seed-to-plant:seed}{बीज} दर्जनों में।
\end{solution}

\begin{solution}{exo:g4:plants-without-seeds:3}
पत्तों की निचली सतह पर: भूरे बिंदुओं की सुथरी क़तारें, हर एक हज़ारों
\omterm{def:g4:plants-without-seeds:spore}{बीजाणुओं} की पुड़िया, जो पकने पर फट जाती है और उन्हें धूल की तरह उड़ने
देती है।
\end{solution}

\begin{solution}{exo:g4:plants-without-seeds:4}
\omterm{prop:g4:plants-without-seeds:damp}{मॉस} \omterm{def:g4:plants-without-seeds:spore}{बीजाणुओं} से जनन करती है, और \omterm{def:g4:plants-without-seeds:spore}{बीजाणुओं} तथा उनकी नाज़ुक नन्ही अवस्थाओं
को नमी चाहिए --- सूखी हवा में वे सूखकर मर जाती हैं। छायादार उत्तरी
दीवार नम रहती है; धूप वाली दक्षिणी दीवार तप जाती है।
\end{solution}

\begin{solution}{exo:g4:plants-without-seeds:5}
जनक \omterm{def:g1:plants-around-us:plant}{पौधे} के किसी टुकड़े से सीधे नया \omterm{def:g1:plants-around-us:plant}{पौधा} बनाना --- \omterm{ex:g4:plants-without-seeds:runner}{भूस्तारी}, \omterm{ex:g4:plants-without-seeds:bulbtuber}{शल्ककंद},
\omterm{ex:g4:plants-without-seeds:bulbtuber}{कंद}, \omterm{met:g4:plants-without-seeds:cutting}{कलम}। सिर्फ़ एक जनक।
\end{solution}

\begin{solution}{exo:g4:plants-without-seeds:6}
\omterm{def:g1:plants-around-us:plant}{पौधा} मिट्टी पर एक लंबा पतला \omterm{def:g1:plants-around-us:parts}{डंठल} --- \omterm{ex:g4:plants-without-seeds:runner}{भूस्तारी} --- फैलाता है; जहाँ वह
ज़मीन छूता है वहाँ \omterm{def:g1:plants-around-us:parts}{जड़ें} जमाता और \omterm{def:g1:plants-around-us:parts}{पत्तियाँ} निकालता है, और एक पूरा नन्हा
\omterm{def:g1:plants-around-us:plant}{पौधा} बन जाता है; फिर \omterm{ex:g4:plants-without-seeds:runner}{भूस्तारी} सूख जाता है और बच्चा अपने जनक जैसा हूबहू,
अपने पैरों पर खड़ा रह जाता है।
\end{solution}

\begin{solution}{exo:g4:plants-without-seeds:7}
बोया हुआ आलू (एक \omterm{ex:g4:plants-without-seeds:bulbtuber}{कंद}) अपनी ``\omterm{def:g1:the-five-senses:sense}{आँखों}'' से भरोसे के साथ \omterm{def:g2:from-seed-to-plant:germination}{अंकुरित होता है},
और हर \omterm{def:g1:the-five-senses:sense}{आँख} से बोई गई अच्छी क़िस्म जैसा हूबहू एक पूरा \omterm{def:g1:plants-around-us:plant}{पौधा} उग सकता है ---
\omterm{def:g2:from-seed-to-plant:seed}{बीजों} से ज़्यादा तेज़ और ज़्यादा पक्का।
\end{solution}

\begin{solution}{exo:g4:plants-without-seeds:8}
उत्तर खुला है। डूबे हुए \omterm{def:g1:plants-around-us:parts}{डंठल} से सफ़ेद \omterm{def:g1:plants-around-us:parts}{जड़ें} निकलती हैं, आम तौर पर एक से
दो हफ़्ते में; और तब \omterm{met:g4:plants-without-seeds:cutting}{कलम} रोपने लायक़ हो जाती है --- एक जनक से बना एक
पूरा नया पुदीना।
\end{solution}

\begin{solution}{exo:g4:plants-without-seeds:9}
\omterm{def:g4:plants-without-seeds:spore}{बीजाणु} नंगा और सस्ता है: बिना \omterm{def:g3:animals-through-seasons:active}{भोजन-भंडार} या नन्हे \omterm{def:g1:plants-around-us:plant}{पौधे} के लगभग हर एक
खो जाता है, इसलिए लाखों बनाने पड़ते हैं। \omterm{def:g2:from-seed-to-plant:seed}{बीज} एक महँगा
उत्तरजीविता-सामान है, इसलिए दर्जन भर अच्छी तरह लैस \omterm{def:g2:from-seed-to-plant:seed}{बीज} ही काफ़ी हैं।
इनमें कॉड मछली (लाखों, बिना देखभाल) और हथिनी (कम, अच्छी तरह सँभाले
हुए) की गूँज है।
\end{solution}

\begin{solution}{exo:g4:plants-without-seeds:10}
स्ट्रॉबेरियाँ \omterm{ex:g4:plants-without-seeds:runner}{भूस्तारी} से आईं --- यानी \omterm{def:g4:plants-without-seeds:vegetative}{कायिक जनन}, एक जनक, हूबहू
नक़लें। सेम \omterm{def:g2:from-seed-to-plant:seed}{बीजों} से आए, जो फूलों ने बनाए थे और जिनमें दो जनकों का
योगदान मिला था (एक \omterm{def:g3:flowers-fruits-seeds:parts}{फूल} से \omterm{def:g3:flowers-fruits-seeds:parts}{पराग}, दूसरे में \omterm{def:g4:animal-reproduction:sexual}{अंडाणु}), इसलिए सेम का हर
\omterm{def:g1:plants-around-us:plant}{पौधा} थोड़ा अलग है।
\end{solution}

\begin{solution}{exo:g4:plants-without-seeds:11}
\omterm{def:g4:plants-without-seeds:vegetative}{कायिक जनन} --- वह शाखा एक प्राकृतिक \omterm{met:g4:plants-without-seeds:cutting}{कलम} है। बाढ़ ने माली की कैंची का
काम किया: उसने शाख़ काटी, उसे ढोया, और गीले कीचड़ में रोप दिया, जो बेंत
के \omterm{def:g1:plants-around-us:parts}{जड़} जमाने की सबसे अच्छी क्यारी है। नदियाँ बेंतों को तोड़ने और रोपने
का काम एक ही झटके में करके अपने किनारों पर उनकी क़तार लगा लेती हैं।
\end{solution}
